% GENERATED FILE. Edit canonical lessons, then run npm run generate:books.
% canonical-source-hash: fnv1a64:50b06deb81826df3
% canonical-lessons: ZH-C352-tonggan, ZH-C352-kending, ZH-C352-renke, ZH-C352-benyi, ZH-C352-biaoshi

\chapter{Same Here}
\label{ch:zh-same-here}
\hlchaptermodality{352}

\begin{chapteropening}
\textbf{By the end of this chapter:} \emph{I can say I feel the same, give my approval and say what I really meant.}

Complete the last lesson of chapter 352: I can say I feel the same, give my approval and say what I really meant.
\end{chapteropening}

\section[\texorpdfstring{\zh{同感} (tónggǎn)}{同感 (tónggǎn)}]{\zh{同感} (tónggǎn) — the same feeling}

\label{lesson:ZH-C352-tonggan}



\begin{quote}
Before the new one: say the Mandarin for an ideal to live up to, then the Mandarin for a person's state of mind.
\end{quote}

\subsection*{You'll want to know: \zh{同感}}
\textbf{\zh{同感}} — \emph{tónggǎn} — \textquotedblleft{}the same feeling\textquotedblright{}.

\begin{quote}
\textbf{\zh{我也有同感}} \emph{wǒ yě yǒu tónggǎn} — I feel the same.
\end{quote}

\begin{cousinweb}
\textbf{\zh{同感}} (\emph{tónggǎn}) is built from \zh{同} 'the same' and \zh{感} 'to feel', both inherited from Middle Chinese, \emph{duwng} and \emph{komX} (Guangyun, in Baxter's transcription), which glosses the first 'equal' and 'together', and the second 'to move'.
\end{cousinweb}

\begin{grammarlens}[title={What you've built}]
One more word for giving reasons.
\end{grammarlens}

\subsection*{Your turn}
\begin{itemize}
\raggedright
  \item \emph{Say it:} \emph{tónggǎn}
  \item \emph{Say it:} \emph{tónggǎn}, once more
  \item \emph{Say it:} \emph{wǒ yě yǒu tónggǎn} — I feel the same.
  \item \emph{Recall:} say the Mandarin for an ideal to live up to, then the Mandarin for a person's state of mind, then say \emph{tónggǎn} again
\end{itemize}

\begin{tcolorbox}[breakable,colback=teal!4,colframe=teal!35!black,title={Before you move on}]
What does \zh{同感} mean? (\textquotedblleft{}The same feeling\textquotedblright{}.) Say it once more.
\end{tcolorbox}
\section[kěndìng]{kěndìng — for sure; without doubt}

\label{lesson:ZH-C352-kending}



\begin{quote}
Before the new one: say the Mandarin for a person's state of mind, then the Mandarin for the same feeling.
\end{quote}

\subsection*{You'll want to know: kěndìng}
\emph{kěndìng} — \textquotedblleft{}for sure; without doubt\textquotedblright{}.

\begin{quote}
\emph{wǒ juéde tā kěndìng huì tóngyì} — I'm sure he'll agree.
\end{quote}

Its characters come later; for now it is a word for the ear.

\begin{cousinweb}
\emph{kěndìng} is built from two parts, both inherited from Middle Chinese, \emph{khongX} and \emph{dengH} (Guangyun, in Baxter's transcription), which glosses the first 'to permit' and the second 'to settle'.
\end{cousinweb}

\begin{grammarlens}[title={What you've built}]
One more word for giving reasons.
\end{grammarlens}

\subsection*{Your turn}
\begin{itemize}
\raggedright
  \item \emph{Say it:} \emph{kěndìng}
  \item \emph{Say it:} \emph{kěndìng}, once more
  \item \emph{Say it:} \emph{wǒ juéde tā kěndìng huì tóngyì} — I'm sure he'll agree.
  \item \emph{Recall:} say the Mandarin for a person's state of mind, then the Mandarin for the same feeling, then say \emph{kěndìng} again
\end{itemize}

\begin{tcolorbox}[breakable,colback=teal!4,colframe=teal!35!black,title={Before you move on}]
What does kěndìng mean? (\textquotedblleft{}For sure; without doubt\textquotedblright{}.) Say it once more.
\end{tcolorbox}
\section[rènkě]{rènkě — to give your approval; to accept as good}

\label{lesson:ZH-C352-renke}



\begin{quote}
Before the new one: say the Mandarin for the same feeling, then the Mandarin for for sure.
\end{quote}

\subsection*{You'll want to know: rènkě}
\emph{rènkě} — \textquotedblleft{}to give your approval; to accept as good\textquotedblright{}.

\begin{quote}
\emph{dàjiā dōu rènkě tā de xiǎngfǎ} — Everyone approved of his idea.
\end{quote}

Its characters come later; for now it is a word for the ear.

\begin{cousinweb}
\emph{rènkě} is built from \emph{rèn} and \emph{kě}, both inherited from Middle Chinese, \emph{nyinH} and \emph{khaX} (Guangyun, in Baxter's transcription), which glosses the first 'to recognize' and the second 'to permit'.
\end{cousinweb}

\begin{grammarlens}[title={What you've built}]
One more word for giving reasons.
\end{grammarlens}

\subsection*{Your turn}
\begin{itemize}
\raggedright
  \item \emph{Say it:} \emph{rènkě}
  \item \emph{Say it:} \emph{rènkě}, once more
  \item \emph{Say it:} \emph{dàjiā dōu rènkě tā de xiǎngfǎ} — Everyone approved of his idea.
  \item \emph{Recall:} say the Mandarin for the same feeling, then the Mandarin for for sure, then say \emph{rènkě} again
\end{itemize}

\begin{tcolorbox}[breakable,colback=teal!4,colframe=teal!35!black,title={Before you move on}]
What does rènkě mean? (\textquotedblleft{}To give your approval; to accept as good\textquotedblright{}.) Say it once more.
\end{tcolorbox}
\section[\texorpdfstring{\zh{本意} (běnyì)}{本意 (běnyì)}]{\zh{本意} (běnyì) — what one really meant; original intention}

\label{lesson:ZH-C352-benyi}



\begin{quote}
Before the new one: say the Mandarin for for sure, then the Mandarin for to give your approval.
\end{quote}

\subsection*{You'll want to know: \zh{本意}}
\textbf{\zh{本意}} — \emph{běnyì} — \textquotedblleft{}what one really meant; original intention\textquotedblright{}.

\begin{quote}
\textbf{\zh{这不是我的本意}} \emph{zhè bú shì wǒ de běnyì} — That's not what I meant.
\end{quote}

\begin{cousinweb}
\textbf{\zh{本意}} (\emph{běnyì}) is built from \zh{本} 'the root' and \zh{意} 'intent', both inherited from Middle Chinese, \emph{pwonX} and \emph{'iH} (Guangyun, in Baxter's transcription), which glosses the first 'the root, as against the branches' and the second 'will'.
\end{cousinweb}

\begin{grammarlens}[title={What you've built}]
One more word for giving reasons.
\end{grammarlens}

\subsection*{Your turn}
\begin{itemize}
\raggedright
  \item \emph{Say it:} \emph{běnyì}
  \item \emph{Say it:} \emph{běnyì}, once more
  \item \emph{Say it:} \emph{zhè bú shì wǒ de běnyì} — That's not what I meant.
  \item \emph{Recall:} say the Mandarin for for sure, then the Mandarin for to give your approval, then say \emph{běnyì} again
\end{itemize}

\begin{tcolorbox}[breakable,colback=teal!4,colframe=teal!35!black,title={Before you move on}]
What does \zh{本意} mean? (\textquotedblleft{}What one really meant; original intention\textquotedblright{}.) Say it once more.
\end{tcolorbox}
\section[biǎoshì]{biǎoshì — to make your view known; to express}

\label{lesson:ZH-C352-biaoshi}



\begin{quote}
Before the new one: say the Mandarin for to give your approval, then the Mandarin for what one really meant.
\end{quote}

\subsection*{You'll want to know: biǎoshì}
\emph{biǎoshì} — \textquotedblleft{}to make your view known; to express\textquotedblright{}.

\begin{quote}
\emph{yīnwèi dàjiā dōu tóngyì suǒyǐ tā yě biǎoshì tóngyì} — Everyone agreed, so he said he agreed too.
\end{quote}

Its characters come later; for now it is a word for the ear.

\begin{cousinweb}
\emph{biǎoshì} is built from \emph{biǎo} and \emph{shì}, both inherited from Middle Chinese, \emph{pjewX} and \emph{zyijH} (Guangyun, in Baxter's transcription), which glosses the first 'to make clear' and the second 'to hand down a sign'.
\end{cousinweb}

\begin{grammarlens}[title={What you've built}]
One more word for giving reasons.
\end{grammarlens}

\subsection*{Your turn}
\begin{itemize}
\raggedright
  \item \emph{Say it:} \emph{biǎoshì}
  \item \emph{Say it:} \emph{biǎoshì}, once more
  \item \emph{Say it:} \emph{yīnwèi dàjiā dōu tóngyì suǒyǐ tā yě biǎoshì tóngyì} — Everyone agreed, so he said he agreed too.
  \item \emph{Recall:} say the Mandarin for to give your approval, then the Mandarin for what one really meant, then say \emph{biǎoshì} again
\end{itemize}

\begin{tcolorbox}[breakable,colback=teal!4,colframe=teal!35!black,title={Before you move on}]
What does biǎoshì mean? (\textquotedblleft{}To make your view known; to express\textquotedblright{}.) Say it once more.
\end{tcolorbox}
